\subsection{Simplicial homotopy}

\bdefn

    Let $f,g\colon K\to X$ be simplicial maps.
    A {\emphcolor simplicial homotopy} from $f$ to $g$ is a map $h\colon K\times\Delta^1\to X$ such that the following diagram
    commutes

    \centerline{\drawdiagram{
        $K\times\Delta^0=K$\cr
        $K\times\Delta^1$&$X$\cr
        $K\times\Delta^0=K$
    }{
        \diagarrow{from={1,1}, to={2,1}, text={1\times\delta^1}, x distance=-.5cm}
        \diagarrow{from={3,1}, to={2,1}, text={1\times\delta^0}, x distance=-.5cm}
        \diagarrow{from={1,1}, to={2,2}, text={f}, x distance=.5cm}
        \diagarrow{from={3,1}, to={2,2}, text={g}, x distance=.25cm}
        \diagarrow{from={2,1}, to={2,2}, text={h}, y distance=.25cm}
    }}

    If there exists a simplicial homotopy from $f$ to $g$ we write $f\hto g$.

\edefn

Recall that a $k$-simplex in $\Delta^n$ is an increasing $f\colon[k]\to[n]$ which we can represent as $[f(0),\dots,f(k)]$.
In particular its vertices can naturally be written as $0,\dots,n$ according to the image of $f\colon[0]\to[n]$.

Note that $\delta^1\colon\Delta^0\to\Delta^1$ corresponds to the vertex of $\Delta^1$ given by $\delta^1\colon[0]\to[1]$,
which by definition skips $1$; its image is $0$.
Similarly $\delta^0$ corresponds to the vertex $1$.

So a homotopy can equivlantely be given as a simplicial map $h\colon K\times\Delta^1\to X$ such that
$$ h(x,0) = f(x),\qquad h(x,1) = g(x) $$
for all $x\in K$ (where $0,1$ denote the degeneracies of the vertices $0,1$).

\bdefn

    Suppose $\iota\colon L\inj K$ is an inclusion of simplicial sets, and let $f,g\colon K\to X$ be simplicial maps.
    Further suppose that $f|_L=g|_L$.
    A homotopy from $f$ to $g$ {\emphcolor relative to $L$} is a homotopy $h\colon f\hto g$ such that the following diagram
    commutes.

    \bigskip
    \centerline{\drawdiagram{
        $K\times\Delta^1$&$X$\cr
        $L\times\Delta^1$&$L$\cr
    }{
        \diagarrow{from={1,1}, to={1,2}, text={h}, y distance=.25cm}
        \diagarrow{from={2,1}, to={1,1}, text={\iota\times\id}, x distance=-.5cm, left cap={(}}
        \diagarrow{from={2,1}, to={2,2}, text={\proj_L}, y distance=-.25cm}
        \diagarrow{from={2,2}, to={1,2}, text={f|_L=g|_L}, x distance=.5cm}
    }}

\edefn

This means that $h(x,t)=f(x)=g(x)$ for all $x\in L,t\in\Delta^1$.
A homotopy of the form
$$ L\times\Delta^1 \xtos{\proj_L} L \xtos\alpha X $$
is called a constant homotopy: it is a homotopy from $\alpha$ to itself.

\bnote

    Homotopy is not in general an equivalence.
    Consider the maps $\iota_0,\iota_1\colon\Delta^0\to\Delta^n$ classifying the vertices $0,1$ respectively.
    A homotopy between them is a map $h\colon\Delta^1\to\Delta^n$ such that $h(0)=0$ and $h(1)=1$.
    This corresponds to the $1$-simplex $[0,1]\in\Delta^n$.

    But there is not a homotopy from $\iota_1$ to $\iota_0$ as this would correspond to a simplex $[1,0]$ which does not
    exist.

\enote

For an $n$-simplex $\sigma\in X$, denote its {\emphcolor boundary} by
$$ \partial\sigma = (d_0\sigma,\dots,d_n\sigma) . $$

\blemm

    Suppose that $X$ is a fibrant simplicial set.
    Then the simplicial homotopy of vertices $x\colon\Delta^0\to X$ is an equivalence.

\elemm

\bproof

    There is a homotopy $x\hto y$ iff there is a $1$-simplex $v\in X_1$ such that $d_1v=x$ and $d_0v=y$.
    That is, iff there is a $1$-simplex $v$ with boundary $\partial v=(y,x$.
    Note that $\partial(s_0x)=(x,x)$ gives reflexivity.

    Now suppose $\partial v_2=(y,x)$ and $\partial v_0=(z,y)$ so that $d_0v_2=d_1v_0=y$.
    Thus $v_2,v_0$ determine a map $(v_0,\ ,v_2)\colon\Lambda^2_1\to X$.
    Because $X$ is fibrant, we can lift it,

    \centerline{\drawdiagram{
        $\Lambda^2_1$&$X$\cr
        $\Delta^2$
    }{
        \diagarrow{from={1,1}, to={2,1}, left cap={(}}
        \diagarrow{from={2,1}, to={1,2}, text={\theta}, y distance=-.25cm}
        \diagarrow{from={1,1}, to={1,2}, text={(v_0,\ ,v_2)}, y distance=.25cm}
    }}

    Then
    $$ \partial(d_1\theta) = (d_0d_1\theta,d_1d_1\theta) = (d_0d_0\theta,d_1d_2\theta) , $$
    and
    $$ d_0\theta = \theta\circ\delta^0 = (v_0,\ ,v_2)\circ\delta^0 = v_0 $$
    and similar for $d_2\theta$.
    Thus
    $$ \partial(d_1\theta) = (d_0v_0,d_1v_2) = (z,x) $$
    so that $d_1\theta$ is a homotopy $x\hto z$; so homotopy is transitive.

    For symmetry, suppose $\partial v_2=(y,x)$ and let $v_1=s_0x$.
    Then $d_1v_1=d_1v_2=x$, so these define a map $(\ ,v_1,v_2)\colon\Lambda^2_0\to X$.
    We extend this to $\theta'\colon\Delta^1\to X$ so that
    $$ \partial(d_0\theta') = (d_0d_0\theta',d_1d_0\theta') = (d_0d_1\theta',d_0d_2\theta') = (x,y) . $$
    Thus the relation is symmetric.
    \qed

\eproof

\blemm[name={fibrantfibers}]

    The pullback of a fibration is a fibration.
    Explicitly, suppose we have the following pullback square

    \centerline{\drawdiagram{
        $F$&$X$\cr
        $W$&$Y$
    }{
        \diagarrow{from={1,1}, to={1,2}}
        \diagarrow{from={1,1}, to={2,1}, text={q}, x distance=-.25cm}
        \diagarrow{from={1,2}, to={2,2}, text={p}, x distance=.25cm}
        \diagarrow{from={2,1}, to={2,2}}
    }}

    If $p$ is a fibration, then so is $q$.

    This implies that all the fibers of $p$ are fibrant simplicial sets.

\elemm

\bproof

    Consider the diagram

    \centerline{\drawdiagram{
        $\Lambda^n_k$&$F$&$X$\cr
        $\Delta^n$&$W$&$Y$\cr
    }{
        \diagarrow{from={1,1}, to={1,2}}
        \diagarrow{from={1,1}, to={2,1}, left cap={(}}
        \diagarrow{from={1,2}, to={2,2}, text={q}, x distance=-.25cm}
        \diagarrow{from={2,1}, to={2,2}}
        \diagarrow{from={1,2}, to={1,3}}
        \diagarrow{from={1,3}, to={2,3}, text={p}, x distance=.25cm}
        \diagarrow{from={2,2}, to={2,3}}
        \diagarrow{from={2,1}, to={1,3}, dashed, text={\theta}, x distance=.25cm}
    }}

    Because $p$ is fibrant, there is a lifting $\theta\colon\Delta^n\to X$.
    So we have maps $\Delta^n\to X,W$ which commute with the pullback diagram, and thus define a unique map $\tilde\theta\colon\Delta^n\to F$.
    This $\tilde\theta$ commutes with $F\xto qW$ and $\Delta^n\to W$ by definition.

    It then follows that it commutes with $\Lambda^n_k\inj\Delta^n$ and $\Lambda^n_k\to F$ by the universal property of the pullback.
    That is, it is easy to see that the composite $\Lambda^n_k\inj\Delta^n\xto{\tilde\theta}F$ composes with $F\to X$ and $F\to W$ the same way that $\Lambda^n_k\to F$
    does.
    Because $F$ is the pullback, this means that the composite is equal to $\Lambda^n_k\to F$ as desired.

    The last part of the claim follows because the fibers of $p$ are the pullbacks along $*\to Y$, and then $q$ is the terminal map $F\to*$.
    It being a fibration means that $F$ is fibrant.
    \qed

\eproof

\bcoro

    If $X$ is fibrant and $L\subseteq K$ is an inclusion, then
    \benum
        \item homotopy of maps $K\to X$ is an equivalence, and
        \item homotopy of maps $K\to X\hrel L$ is an equivalence.
    \eenum

\ecoro

\bproof

    Clearly $(2)$ implies $(1)$ with $L=\emptyset$.
    Now we know that $i^*\colon\ihom(L,X)\to\ihom(K,X)$ is a fibration, and the fiber of a vertex $\alpha\in\ihom(K,X)_0$
    consists of all $f\in\ihom(L,X)$ such that $f|_L=\alpha$ (recall $\ihom(K,X)_0=\hom_\sSet(k,X)$).
    These fibers are all fibrant by the above lemma.

    Now we claim that $f\hto g\hrel L$ iff $f,g$ are homotopic as vertices of the fiber of $i^*$.
    Let $\alpha=f|_L=g|_L$.
    A homotopy $h\colon h\hto g\hrel L$ corresponds to a map $h\colon\Delta^1\times K\to X$, which by the \refmath{explaw}
    corresponds to $\tilde h\colon\Delta^1\to\ihom(K,X)$.

    Note that a homotopy diagram of the form

    \centerline{\drawdiagram{
        $K\times\Delta^0=K$\cr
        $K\times\Delta^1$&$X$\cr
        $K\times\Delta^0=K$
    }{
        \diagarrow{from={1,1}, to={2,1}, text={1\times\delta^1}, x distance=-.5cm}
        \diagarrow{from={3,1}, to={2,1}, text={1\times\delta^0}, x distance=-.5cm}
        \diagarrow{from={1,1}, to={2,2}, text={f}, x distance=.5cm}
        \diagarrow{from={3,1}, to={2,2}, text={g}, x distance=.25cm}
        \diagarrow{from={2,1}, to={2,2}, text={h}, y distance=.25cm}
    }}

    Corresponds to a homotopy of vertices

    \centerline{\drawdiagram{
        $\Delta^0$\cr
        $\Delta^1$&$\ihom(K,X)$\cr
        $\Delta^0$
    }{
        \diagarrow{from={1,1}, to={2,1}, text={\delta^1}, x distance=-.5cm}
        \diagarrow{from={3,1}, to={2,1}, text={\delta^0}, x distance=-.5cm}
        \diagarrow{from={1,1}, to={2,2}, text={\iota_f}, x distance=.5cm}
        \diagarrow{from={3,1}, to={2,2}, text={\iota_g}, x distance=.25cm}
        \diagarrow{from={2,1}, to={2,2}, text={\tilde h}, y distance=.25cm}
    }}

    and vice versa.

    Homotopy relative to $L$ is the same as homotopy of vertices in the fibers of $i^*$, which are fibrant.
    By the above lemma this is an equivalence.
    \qed

\eproof

